\documentclass[9pt,twocolumn,twoside]{pnas-new}
\graphicspath{{figures_and_tables/v006/}{../LatentIndustryRepresentations/results/paper/}}

\articletype{COMPUTER SCIENCES / ECONOMIC SCIENCES}

\templatetype{pnasresearcharticle}
\usepackage{multirow}

\begin{document}

\title{Multi-View Representation Learning for Regional Opportunity Spaces and Industrial Path Development}

\author[a,b,1]{Philipp Stark}
\author[c]{Alexandros Sopasakis}
\author[a]{Ola Hall}
\author[a,b]{Markus Grillitsch}

\affil[a]{Department of Human Geography, Lund University, 221 00 Lund, Sweden}
\affil[b]{CIRCLE (Centre for Innovation Research), Lund University, 221 00 Lund, Sweden}
\affil[c]{Centre for Mathematical Sciences, Lund University, 221 00 Lund, Sweden}

\leadauthor{Stark}

\significancestatement{Environmental pressures, geopolitical tensions, and technological disruption fundamentally change the conditions for industrial development and sustained wellbeing. Understanding industrial transformation is difficult because it depends on complex interactions between industries through multiple types of relations. These form opportunity spaces that vary between places and over time. We develop a methodology that trains machine-learning models on heterogeneous data sources to capture these interrelations and represent evolving regional opportunity spaces for industrial transformation. Using Sweden’s rich, linkable longitudinal data as a testbed, we show that regional opportunity spaces are strong predictors of industrial change over three decades. The methodology works with different data sources and scales, holding strong promise for studying opportunities and limitations for industrial transformation worldwide.}

\authorcontributions{All authors conceived the study. P.S. and M.G. developed the conceptual framing. P.S. led the development of the machine-learning methodology in collaboration with A.S. P.S. collected and processed the data, implemented the computational methods, and performed the analyses. P.S. and M.G. designed the regional evaluation. All authors contributed to the interpretation of the results and the development of the manuscript. A.S., O.H., and M.G. supervised the research. P.S. drafted the manuscript, and all authors critically revised and approved the final version.}
\authordeclaration{The authors declare no competing interests.}
\correspondingauthor{\textsuperscript{1}To whom correspondence should be addressed. E-mail: philipp.stark\@keg.lu.se}

\keywords{regional opportunity spaces $|$ industrial path development $|$ industry relatedness $|$ multi-view representation learning $|$ regional diversification}

\begin{abstract}
Regional opportunities for industrial development arise when local conditions favor particular economic activities. These conditions depend in part on a region's existing industry profile and on the economic relationships among its activities. Heterogeneous data sources on mobility, production, and firm ownership provide complementary evidence about this underlying industrial structure. We therefore develop a multi-view representation-learning framework for industries that integrates these different sources into a shared latent representation. By using graph diffusion, the framework captures both direct and indirect connections between industries. Applied to Swedish industries over three decades and three classification versions, the learned representation places industries from different classification periods within a common latent space. To evaluate this operationalization of opportunity, we construct an opportunity-density measure by combining latent similarity with local industry size. Higher opportunity density is consistently associated with industry entry and startup formation and with lower industry exit. Across classification periods, a one-standard-deviation increase in opportunity density is associated with 5.3--6.2 percentage-point higher industry-entry probabilities (118--136\% of baseline rates), 8.9--13.8 percentage-point lower industry-exit probabilities (147--209\%), and 1.1--1.7 percentage-point higher startup-formation probabilities (131--187\%). Opportunity density remains associated with these outcomes in 41 of 42 models after accounting for information contained in the individual data sources. Our machine-learning framework shows how heterogeneous data can be integrated into a common latent structure and used to operationalize regional opportunity spaces, providing a new approach to studying industrial path development.
\end{abstract}

\dates{This manuscript was compiled on \today}
\doi{\url{www.pnas.org/cgi/doi/10.1073/pnas.XXXXXXXXXX}}


\maketitle
\thispagestyle{firststyle}
\ifthenelse{\boolean{shortarticle}}{\ifthenelse{\boolean{singlecolumn}}{\abscontentformatted}{\abscontent}}{}

\Firstpage

%The \Firstpage command is used to format the first page text column size. The same size will be maintained for subsequent paragraph until the \Endparasplit or \Parasplit command is encountered.
\section{Introduction}
Regional opportunities for industrial development arise when regional conditions favor some economic activities over others. The concept of a regional opportunity space captures these conditions that make some activities more favourable than others at a given time. These possibilities depend on the capabilities and constraints present in a region, but also on actors that recognize, pursue, and reshape them. Opportunity spaces therefore vary across regions, over time, and between actors, and include both conditions that enable development and constraints that may prevent it. Understanding regional development consequently requires understanding the opportunity space within which new economic activities can emerge. \cite{grillitschTrinityChangeAgency2020,kurikkaRegionalOpportunitySpaces2023,grillitschPatternsOpportunitySpaces2025,boschmaRelatednessDriverRegional2017}

However, regional opportunity spaces are latent constructs and cannot be observed directly. Their different components must instead be inferred from observable evidence about regional conditions ~\cite{hidalgoBuildingBlocksEconomic2009}. One important component is the industrial structure of a region: the industries present, the relations among them, and the economic structures through which they are connected. \Parasplit Heterogeneous data sources provide partial evidence about this industrial-structural component, but no single source captures it in its entirety~\cite{neffkeSkillRelatednessFirm2013,diodatoWhyIndustriesCoagglomerate2018,hobergTextBasedNetworkIndustries2016}.

Moreover, these sources provide different views on the same phenomena. They may share common patterns while also capturing distinct and complementary relations between industries~\cite{diodatoWhyIndustriesCoagglomerate2018,liEvaluatingPrincipleRelatedness2024,kivelaMultilayerNetworks2014}. Representing the industrial-structural component of regional opportunity spaces therefore requires an approach that can integrate heterogeneous relational evidence while retaining the information contained in its different views.

Several research traditions infer regional development from observable economic structure.
The principle of relatedness shows that regions are more likely to develop activities related
to those already present~\cite{hidalgoPrincipleRelatedness2018,hidalgoEconomicComplexityTheory2021},
while research on industrial branching links such development to capabilities associated with
existing activities~\cite{boschmaRelatednessDriverRegional2017,neffkeHowRegionsDiversify2011}.
Product-space and economic-complexity approaches similarly infer development possibilities
and underlying productive capacities from observed activity structures
~\cite{hidalgoProductSpaceConditions2007,hidalgoBuildingBlocksEconomic2009,
ballandSmartSpecializationPolicy2019}.

\begin{figure*}[t]
    \centering
    \includegraphics[width=\textwidth]
    {figures_and_tables/model_architecture_overview.pdf}
    \caption{\textbf{Overview of multi-view representation learning.}
    Official industry descriptions are encoded as 384-dimensional sentence
    embeddings and reduced by principal component analysis to initialize
    64-dimensional latent industry vectors; text is used only for
    initialization. Labour mobility, common ownership, input--output relatedness,
    the industry hierarchy, and classification concordances are converted into
    fixed training targets at one or more diffusion scales, with a direct
    ordinal target for the hierarchy. The shared latent vectors define one
    cosine-similarity geometry. For each categorical target,
    temperature-scaled cosine scores are converted to a predicted row
    distribution and compared with the empirical target using
    Kullback--Leibler divergence or cross-entropy; the direct hierarchy target
    uses an ordinal margin loss. The weighted losses jointly update the latent
    vectors and target temperatures. The normalized learned vectors define the
    latent similarity matrix used in subsequent analyses.}
    \label{fig:model-architecture}
\end{figure*}

Yet regional development is not restricted to a single form of relatedness. Different economic relations can reveal different aspects of industrial structure: worker movements provide evidence about shared skill requirements~\cite{neffkeSkillRelatednessFirm2013}, while production linkages and co-agglomeration patterns capture dependencies and complementarities between industries~\cite{diodatoWhyIndustriesCoagglomerate2018,delgadoClustersConvergenceEconomic2014}. Industries combined within the same enterprise provide another form of relational evidence, reflecting organizationally feasible combinations of activities and potentially shared resources and economies of scope~\cite{teeceEconomiesScopeScope1980,bryceGeneralInterindustryRelatedness2009}. Similarity in the activities performed by industries, as encoded in official industry descriptions, provides yet another perspective on how industries are connected~\cite{hobergTextBasedNetworkIndustries2016}. Development can also involve combinations beyond established patterns of relatedness~\cite{frenkenRelatedVarietyUnrelated2007,grillitschUnrelatedKnowledgeCombinations2018}.

What is missing is a way to represent these heterogeneous relational structures jointly, preserving their information richness while balancing shared and complementary information. To address this, we develop a self-supervised multi-view representation-learning framework, building on graph and multi-view representation learning~\cite{hamiltonRepresentationLearningGraphs2017,parkRelationalKnowledgeDistillation2019,dongMetapath2vecScalableRepresentation2017,shiMvn2vecPreservationCollaboration2018,schlichtkrullModelingRelationalData2018}. The framework represents the different relational sources (worker movements, production linkages, etc.) as distinct graphs over the same set of industry nodes, which we refer to as relational views in the following. Additionally, industry descriptions (labels describing the industry) provide an initial representation of similarity between industrial activities and are encoded once using a small pretrained sentence encoder. On top of this initialization, each relational view captures a different aspect of the economic structure connecting industries, and these signals are used to train the embedding. Because relevant connections may extend beyond directly linked industry pairs, graph diffusion is used to capture both direct and indirect relational structure within each view~\cite{coifmanDiffusionMaps2006,gasteigerDiffusionImprovesGraph2019,bertagnolliDiffusionGeometryMultiplex2021,zhangGraphDiffusionEnhanced2026}. This means that indirect (second-, third-order, etc.) relations of an industry are also mined from the graphs and learned. Together with industry hierarchy and classification concordance constraints, these views jointly shape a shared latent representation in which each industry is represented by a 64-dimensional embedding vector. Similarity between these vectors defines latent proximity between industries and provides the basis for the regional evaluation that follows (Figure~\ref{fig:model-architecture}).

To operationalise regional opportunity spaces, we combine cosine similarity in the latent industry representation with local industry size (number of employees) to construct an opportunity-density measure. Opportunity density captures how closely a focal industry relates to a region's existing industry portfolio, with industries weighted by their local employment. Higher opportunity density therefore means that the region's employment-weighted industry portfolio is, on average, more closely related to the focal industry. We interpret such a configuration as indicating that the existing industrial structure provides a combination of conditions that is favourable to the development of activities in the focal industry. Opportunity density is thus used to evaluate regional opportunity spaces for industrial path development.

We apply the framework to Swedish industries over three decades, spanning three successive versions of the Swedish Standard Industrial Classification (SNI92, SNI02, and SNI07), Sweden's implementation of NACE~\cite{statisticsswedenscbMIS200722007,europeancommissionNACERev22008}. We evaluate whether the learned representation reconstructs the underlying relational views, preserves meaningful industry structure, and links industries across successive classification systems. For the regional evaluation, we examine whether opportunity density is associated with subsequent industry entry, industry exit, and startup formation, and whether these associations remain when densities derived from the individual relational views are included in the same models.

\section{Results}
\label{sec:results}
We first examine how the relational views differ and how well the shared latent representation captures their relationships. We then assess its classification structure and cross-version correspondence. Finally, we perform the regional evaluation based on the opportunity density measure. We assess whether the associations between opportunity density and regional outcomes persist in held-out municipalities (five-fold spatial cross-fitting) and held-out subsequent years. We also assess their stability across representations trained from perturbed initializations. Finally, we test whether opportunity density carries information beyond the individual relational views.

\subsection{Complementarity of relational views}
\label{sec:results-complementarity}

The relational views capture different aspects of industrial relatedness (Figure~\ref{fig:source-complementarity}). At the common level of SNI07 two-digit divisions, pairwise rank correlations range from $0.02$ between input--output relatedness and the industry hierarchy to $0.51$ between labour mobility and common ownership. The views also differ in their local neighbourhoods: across all pairs of views, the mean Jaccard overlap among their ten nearest neighbours is only $0.16$--$0.26$. Mutual information, which also captures nonlinear dependence, ranges from $0.03$ to $0.14$.

At the broader matrix level, centered kernel alignment is higher among labour mobility, common ownership, and input--output relatedness ($0.89$--$0.95$). Missing pairs are zero-filled for this calculation, so shared sparsity can contribute to the high alignment. The views therefore combine some broad structural agreement with substantial differences in which industries they connect most closely, motivating their joint use.

\begin{figure*}[t]
    \centering
    \includegraphics[width=\textwidth]
    {figures_and_tables/result1_source_complementarity.png}
    \caption{\textbf{Complementarity of relational views.}
    All comparisons are calculated at the SNI07 two-digit division level.
    Labels labor, ownership, io, and hierarchy denote labour mobility,
    common ownership, input--output relatedness, and industry hierarchy.
    Panel A reports pairwise Spearman rank correlations; Panel B reports the
    mean Jaccard overlap between each view's ten nearest neighbours;
    Panel C reports centered kernel alignment; and Panel D reports symmetric
    $k$-nearest-neighbour mutual information. The figure is restricted to the
    relational views included in the reported representation. Diagonal
    values are mechanical; mutual-information diagonal entries are fixed at one rather than estimated self-information.}
    \label{fig:source-complementarity}
\end{figure*}


\subsection{Reconstructing heterogeneous evidence in one geometry}
\label{sec:results-reconstruction}

The shared latent representation more closely matches the economic relations than the original text encoder does (Table~\ref{tab:rho-table}). Mean Spearman correlations rise from $0.422$ to $0.595$ for labour mobility and from $0.449$ to $0.712$ for common ownership. For input--output relatedness, the correlation reaches $0.911$ at the two-digit level, compared with $0.180$ for the text encoder; averaging four-digit cosine similarities to division pairs gives $0.665$, compared with $0.219$.

This comparison uses the original 384-dimensional text encoder, before its reduction to the 64-dimensional semantic initialization. It therefore compares two complete representations, rather than isolating training at a fixed dimension. Greater alignment with economic relations comes with lower alignment with the industry hierarchy: the Spearman correlation between cosine distance and classification-tree distance, computed on pairs balanced across the four hierarchy tiers, falls from $0.678$ to $0.372$ (Materials and Methods). The shared latent geometry thus accommodates several economic relationships while retaining a weaker correspondence with the classification tree. Period-specific reconstruction correlations and evaluation KL are reported in SI Appendix, Table~S1.

\begin{table}[t]
\centering
\begin{tabular}{lrrr}
\toprule
Relational view & Text encoder & Joint & $\Delta$ \\
\midrule
Labour mobility & 0.4221 & 0.5946 & 0.1726 \\
Common ownership & 0.4491 & 0.7115 & 0.2624 \\
Input--output (2-digit) & 0.1804 & 0.9110 & 0.7306 \\
Input--output (4-digit) & 0.2191 & 0.6652 & 0.4461 \\
Industry hierarchy & 0.6779 & 0.3721 & -0.3058 \\
\bottomrule
\end{tabular}
\caption{Reconstruction of relational views: original text encoder (384 dimensions) versus joint representation (64 dimensions). Labour mobility and common ownership report mean Spearman correlations across classification periods. Input--output correlations use SNI07 two-digit hub cosines or four-digit cosines averaged to division pairs. Industry hierarchy compares cosine distance with tree distance on pairs balanced across hierarchy tiers: all same-group and same-division pairs, plus 5,000 sampled pairs from each of the same-section and different-section tiers. $\Delta$ is joint minus text encoder.}
\label{tab:rho-table}
\end{table}


\subsection{Structure of the shared latent representation}
\label{sec:results-emergent}

\subsubsection{Hierarchy preservation}

The shared latent geometry retains a graded relationship with the industry hierarchy (Figure~\ref{fig:hierarchy-preservation}). Mean cosine similarity between SNI07 four-digit industries is $0.529$ within the same three-digit group, $0.381$ within the same two-digit division but different groups, $0.177$ within the same section but different divisions, and $0.140$ between different sections. Although the distributions overlap, this ordering shows that classificatory proximity remains visible alongside the economic relationships incorporated during training.

\begin{figure}[t]
    \centering
    \includegraphics[width=\columnwidth]{figures_and_tables/result3_hierarchy_preservation.pdf}
    \caption{\textbf{Preservation of the official industry hierarchy.}
    Cosine similarity between all 188,805 pairs of SNI07 four-digit industries in the learned representation, by the closest shared level of the classification. Boxes show quartiles, vertical lines medians, diamonds means, and whiskers the 5th to 95th percentiles. Mean similarities and numbers of pairs: $0.529$ (865 pairs), $0.381$ (3,594), $0.177$ (29,012), and $0.140$ (155,334).}
    \label{fig:hierarchy-preservation}
\end{figure}


\subsubsection{Cross-version retrieval}

The shared latent representation also links industries across classification versions. For SNI02-to-SNI07, a valid successor is the nearest neighbour for $71.8\%$ of codes and appears among the ten nearest neighbours for $97.7\%$. The corresponding rates for SNI92-to-SNI02 are $88.3\%$ and $99.8\%$.

The original 384-dimensional text encoder performs better on these retrieval diagnostics (Table~\ref{tab:retrieval-table}). Mean average precision falls from $0.733$ to $0.700$ for SNI02-to-SNI07 and from $0.976$ to $0.926$ for SNI92-to-SNI02. The joint representation therefore retains substantial cross-version correspondence while accommodating the additional economic relations. Because classification concordances also guide training, these results describe alignment among known codes rather than retrieval for unseen industries.

\begin{table}[t]
\centering
\begin{tabular}{lrrrr}
\toprule
Direction & Metric & Text encoder & Joint & $\Delta$ \\
\midrule
SNI02 $\rightarrow$ SNI07 & MRR & 0.9322 & 0.8126 & -0.1196 \\
SNI02 $\rightarrow$ SNI07 & MAP & 0.7334 & 0.7003 & -0.0331 \\
SNI02 $\rightarrow$ SNI07 & Hit@1 & 0.9047 & 0.7179 & -0.1868 \\
SNI02 $\rightarrow$ SNI07 & Hit@10 & 0.9805 & 0.9767 & -0.0038 \\
SNI92 $\rightarrow$ SNI02 & MRR & 0.9920 & 0.9333 & -0.0587 \\
SNI92 $\rightarrow$ SNI02 & MAP & 0.9755 & 0.9256 & -0.0499 \\
SNI92 $\rightarrow$ SNI02 & Hit@1 & 0.9901 & 0.8832 & -0.1069 \\
\bottomrule
\end{tabular}
\caption{Cross-version retrieval: original text encoder (384 dimensions) versus joint representation (64 dimensions). MRR is mean reciprocal rank of the first valid successor; MAP is mean average precision over valid successors; Hit@$k$ is the fraction of source codes with at least one valid successor among the first $k$ retrieved codes. $\Delta$ is the joint minus text encoder.}
\label{tab:retrieval-table}
\end{table}


\subsubsection{The learned industry space}

A two-dimensional t-SNE projection of the 615 SNI07 four-digit industries shows one organization of the shared latent representation (Figure~\ref{fig:learned-opportunity-space}). Colours indicate SNI07 sections. The example neighbourhoods marked in the figure each combine industries from several sections. In the original latent geometry, the SNI07 section silhouette score is $-0.0461$, indicating that sections overlap rather than forming distinct clusters.

\begin{figure*}[t]
    \centering
    \includegraphics[width=\textwidth]
    {figures_and_tables/result3_opportunity_space_new.pdf}
    \caption{\textbf{The shared latent representation of industries.} Two-dimensional t-SNE projection of the embeddings of all 615 SNI07 industries (4-digit codes); colours denote SNI07 sections. Dashed frames mark example neighbourhoods. Each neighbourhood combines industries from several sections. Industry labels are abbreviated. The projection is used only for visualization; all quantitative analyses use the original latent geometry.}
    \label{fig:learned-opportunity-space}
\end{figure*}


\subsection{Opportunity density and regional development outcomes}
\label{sec:results-opportunity-spaces}

We next assess whether the learned geometry captures industrial conditions associated with regional development. Opportunity density measures a focal industry's cosine similarity to a municipality's existing industries, weighted by their log employment and excluding the focal industry itself. Its profile across industries describes the municipality's regional opportunity structure.

We relate density to subsequent industry entry, industry exit, and startup formation, and in SI Appendix to positive employment and wage growth. Estimates are calculated separately for the three classification periods and expressed per one-standard-deviation increase in density. The main models include municipality$\times$year and industry fixed effects, so associations compare industries within the same municipality and year, after accounting for persistent industry differences. Standard errors are clustered by municipality and industry (Materials and Methods). Binary-outcome coefficients are linear associations in percentage points, rather than causal effects or changes in predicted probabilities for individual observations. SI Appendix, Table~S3 reports sample sizes, baseline rates, and period estimates.

\subsubsection{Industry entry}

Higher opportunity density is associated with industry entry in all three classification periods (Figure~\ref{fig:industry-entry}). Observed entry rates rise across density deciles. In the adjusted models, a one-standard-deviation increase in density is associated with an annual entry probability higher by $6.05$ percentage points during SNI92 (95\% CI $4.48$--$7.62$), $6.16$ during SNI02 ($3.93$--$8.40$), and $5.30$ during SNI07 ($3.53$--$7.08$; all $p<0.001$). For context, baseline entry rates are $5.13\%$, $5.10\%$, and $3.89\%$, respectively. Year-specific coefficients are positive in every base year. The association therefore recurs across periods and years, linking the appearance of meaningful local activity to the existing regional industry portfolio.

\begin{figure*}[t]
    \centering
    \includegraphics[width=\textwidth]{figures_and_tables/result4_entry.png}
\caption{\textbf{Opportunity density and regional industry entry.}
Panel A shows mean entry probability (0--1 scale) by within-period
density decile; coloured horizontal dashed lines show baseline rates.
Panel B shows year-specific coefficients with 95\% confidence intervals;
dotted vertical lines mark classification changes. Panel C shows period
coefficients with 95\% confidence intervals. Annotations give effects
in percentage points and as percentages of the baseline rate. Axes
labelled ``Effect'' and ``in \%'' report percentage-point associations
per one-standard-deviation increase in density, using the standard
deviation pooled across periods. Candidate-panel sizes are
$N=988\,126$ (SNI92), $N=353\,657$ (SNI02), and $N=1\,266\,700$
(SNI07) municipality--industry--year observations. Models include
municipality$\times$year and industry fixed effects; year-specific
coefficients come from density$\times$base-year interactions in the
same model. Standard errors are clustered by municipality and industry.}
    \label{fig:industry-entry}
\end{figure*}


\subsubsection{Industry exit}

Higher opportunity density is also associated with lower industry exit rates (Figure~\ref{fig:industry-exit}). The adjusted annual exit probability is lower by $13.77$ percentage points per standard deviation of density in SNI92 (95\% CI $11.69$--$15.85$), $8.91$ in SNI02 ($7.37$--$10.44$), and $11.84$ in SNI07 ($10.18$--$13.49$; all $p<0.001$), against baseline exit rates of $8.07\%$, $6.08\%$, and $5.65\%$. Year-specific coefficients are negative in every base year. Together, the entry and exit results link opportunity density both to the emergence of local industrial activity and to its continued presence.

\begin{figure*}[t]
    \centering
    \includegraphics[width=\textwidth]{figures_and_tables/result4_exit.png}
\caption{\textbf{Opportunity density and regional industry exit.}
Panel A shows mean exit probability (0--1 scale) by within-period
density decile; coloured horizontal dashed lines show baseline rates.
Panel B shows year-specific coefficients with 95\% confidence intervals;
dotted vertical lines mark classification changes. Panel C shows period
coefficients with 95\% confidence intervals. Annotations give effects
in percentage points and as percentages of the baseline rate. Axes
labelled ``Effect'' and ``in \%'' report percentage-point associations
per one-standard-deviation increase in density, using the standard
deviation pooled across periods. Candidate-panel sizes are
$N=622\,823$ (SNI92), $N=242\,582$ (SNI02), and $N=873\,496$
(SNI07) municipality--industry--year observations. The SNI92 sample
contains 289 municipalities; the other samples contain 290. Models
include municipality$\times$year and industry fixed effects;
year-specific coefficients come from density$\times$base-year
interactions in the same model. Standard errors are clustered by
municipality and industry.}
    \label{fig:industry-exit}
\end{figure*}

\subsubsection{Startup formation}
\label{sec:results-start-up}

The association also extends to startup formation: the founding of new, single-establishment firms with at most ten employees, excluding mergers and continuations (Materials and Methods). This outcome occurs in only $0.71\%$--$0.92\%$ of candidate municipality--industry--year observations, making it considerably rarer than industry entry.

A one-standard-deviation increase in opportunity density is associated with an annual probability of startup formation higher by $1.05$ percentage points in SNI92 (95\% CI $0.82$--$1.29$), $1.17$ in SNI02 ($0.54$--$1.80$), and $1.72$ in SNI07 ($1.14$--$2.30$; all $p<0.001$; Figure~\ref{fig:startup-entry}). Year-specific coefficients are positive in every base year. Opportunity density is thus informative about new firm founding as well as the appearance and continued presence of industry-level employment.

\begin{figure*}[t]
    \centering
    \includegraphics[width=\textwidth]{figures_and_tables/result4_startup.png}
\caption{\textbf{Opportunity density and startup formation.}
Panel A shows mean formation probability (0--1 scale) by within-period
density decile; coloured horizontal dashed lines show baseline rates.
Panels B and C show year-specific and period coefficients, respectively,
with 95\% confidence intervals. Effects are percentage points per
one-standard-deviation increase in density, using the standard
deviation pooled across periods; dotted vertical lines mark
classification changes. Panel C annotations give effects in percentage
points and as percentages of baseline rates. Candidate-panel sizes are
$N=988\,126$ (SNI92), $N=353\,657$ (SNI02), and $N=1\,266\,700$
(SNI07) municipality--industry--year observations. Startups are newly
founded single-establishment firms with at most ten employees
(Materials and Methods). Models include municipality$\times$year and
industry fixed effects; year-specific coefficients come from
density$\times$base-year interactions in the same model. Standard
errors are clustered by municipality and industry. The smaller
coefficient for base year 1993 may reflect incomplete firm-demography
coding in outcome year 1994. The plot's ``start-up'' labels refer to
startup formation; ``Effect'' and ``in \%'' denote percentage-point
associations.}
    \label{fig:startup-entry}
\end{figure*}


\subsubsection{Employment and wage growth}
\label{sec:results-growth}

We also examined whether opportunity density is associated with positive employment and wage growth. These associations are small relative to baseline rates ($4$--$21\%$) and do not remain positive and distinguishable from zero once labour-mobility density is included; the SNI07 wage-growth association reverses sign (SI Appendix, Section~2.4). We therefore focus on industry entry, industry exit, and startup formation.

\subsubsection{Spatial and temporal hold-out}
\label{sec:results-holdout}

The relational views used in training and the regional outcomes partly derive from the same register data. We therefore re-estimate the main models with representations retrained without the municipalities or years in which density is evaluated (Materials and Methods; SI Appendix, Tables~S4--S6).

In the spatial hold-out, municipalities are divided into five folds, and the density of each municipality is calculated from a representation trained without the labour mobility and common ownership observed in its fold. The associations remain of similar size: $4.04$--$5.44$ percentage points for industry entry ($76$--$90\%$ of the main estimates), $-13.95$ to $-9.15$ for industry exit ($101$--$103\%$), and $0.97$--$1.40$ for startup formation ($82$--$92\%$; all $p<0.01$). Across three random assignments of municipalities to folds, the coefficients differ by at most $0.3$ percentage points.

In the temporal hold-out, representations are trained only on register data up to the year preceding the explanatory year, and densities for 2016--2018 are related to outcomes in 2017--2019. A one-standard-deviation increase in density is associated with an entry probability higher by $4.06$ percentage points (95\% CI $2.20$--$5.92$), an exit probability lower by $10.92$ ($9.30$--$12.55$), and a startup-formation probability higher by $1.37$ ($0.67$--$2.06$). These coefficients are $77\%$, $97\%$, and $75\%$ of those of the main representation on the same observations. Across ten representations trained from perturbed initializations, the standard deviations of the coefficients are at most $0.17$ percentage points, less than a fifth of their standard errors.

\subsubsection{Opportunity density and the individual relational views}

To assess whether the joint representation carries information beyond its individual views, we estimate multiple regressions that include opportunity density together with the density derived from one relational view, or from the untrained text initialization, using the same fixed effects and candidate panels (Figure~\ref{fig:robustness}; SI Appendix, Table~S7). Opportunity density has a 95\% confidence interval excluding zero in 41 of the 42 models. Densities based on input--output relatedness and the text initialization have intervals including zero once opportunity density is included, while opportunity density retains coefficients of $3.95$--$5.72$ for entry, $-11.15$ to $-5.35$ for exit, and $0.70$--$1.74$ percentage points for startup formation. Common ownership and the industry hierarchy retain smaller associations of their own with industry entry ($0.47$--$1.37$ and $0.64$--$1.08$ percentage points). Hierarchy density is also associated with lower exit in all three periods and common-ownership density in SNI92 only; each is associated with startup formation in two of the three periods. In these models, opportunity density retains coefficients of $3.04$--$5.24$ for entry, $-10.49$ to $-5.08$ for exit, and $0.48$--$1.67$ for startup formation.

Labour mobility is the strongest complementary view. For industry entry, the coefficients are similar ($1.80$--$2.71$ for opportunity density and $1.82$--$2.29$ for labour-mobility density), although the opportunity-density interval includes zero in SNI07. For industry exit and startup formation, both densities remain associated with the outcomes. Opportunity density has the larger coefficient magnitudes for exit ($-7.71$ to $-4.76$ versus $-2.18$ to $-1.35$) and startup formation ($0.54$--$1.01$ versus $0.21$--$0.39$). Each density is scaled by its own within-period standard deviation, so these coefficient comparisons describe associations on their respective scales.

\begin{figure*}[t]
    \centering
    \includegraphics[width=\textwidth]{figures_and_tables/result4_multiple_regression.pdf}
\caption{\textbf{Opportunity density and relational-view densities in
multiple regressions.} Each row and period shows one linear probability
model with opportunity density (filled circle) and one relational-view
density or the untrained text initialization (open square). The 42
models cover industry entry, industry exit, and startup formation.
Coefficients are percentage points per standard deviation of each
density, standardized within the classification period; bars show
95\% confidence intervals. The candidate-panel sizes for SNI92,
SNI02, and SNI07, respectively, are $N=988\,126$, $353\,657$, and
$1\,266\,700$ for entry; $N=622\,823$, $242\,582$, and $873\,496$
for exit; and $N=988\,126$, $353\,657$, and $1\,266\,700$ for
startup formation. Here $N$ counts municipality--industry--year
observations used for each outcome-period comparison, not independent
model estimates. Models include municipality$\times$year and industry
fixed effects; standard errors are clustered by municipality and
industry. Colours mark classification periods; input--output
relatedness is unavailable for SNI92. Estimates are reported in
SI Appendix, Table~S7.}
    \label{fig:robustness}
\end{figure*}


\section{Discussion}
In this study, we combine different sources of information about how industries are related to learn a shared industry space. By taking regional employment profiles into account, we use this space to evaluate the industrial-structural component of regional opportunity spaces. A central result is that the learned industry space is neither a reproduction of any single source nor a simple average of them. Rather, it captures a distinct organization of industries that cannot be reduced to any one of the underlying sources. The relational views derived from these different sources share some broad structure, but their relatively low rank correlations, mutual information, and nearest-neighbour overlap show that they organize individual industries differently. During learning, the shared representation becomes substantially more aligned with labour mobility, common ownership, and input--output relatedness, while at the same time becoming less closely aligned with the formal industry hierarchy. The resulting geometry therefore reflects a balance between shared and complementary information contained in the different relational views.

This balancing is also visible in the way the learned space relates to the industry classification. The hierarchy remains clearly present: mean cosine similarity declines from industries in the same three-digit group, through different groups in the same division and different divisions in the same section, to industries in different sections. At the same time, the weaker correspondence between cosine distance and classification-tree distance shows that the additional economic information repositions industries relative to a purely classificatory representation. Sections therefore remain visible as broad structure without forming isolated clusters. This is an important property of the shared representation: it preserves meaningful classificatory organization while allowing economic relations from other sources to alter which industries are positioned close to one another.

A further feature of the framework is that the learned industry space incorporates not only direct relations between industries but also higher-order relational structure identified through graph diffusion. By propagating relations through intermediate industries, diffusion captures second- and higher-order connections that are not visible from direct links alone. This is potentially important for understanding the role of intermediary industries that connect otherwise more weakly related parts of the industry structure. Such positions may be relevant to gatekeeping or brokerage functions, through which industries mediate access to different pools of skills, knowledge, production relations, or organizational resources~\cite{broekelCriticalLinksKnowledge2018}. The present framework does not identify such gatekeeping mechanisms directly, but the higher-order relational structure it captures provides a basis for examining them in future work.

The classification concordances provide a further benefit by allowing industries from SNI92, SNI02, and SNI07 to coexist within the same latent space. Despite the changes in classification over the three decades studied, valid successor industries appear among the ten nearest neighbours for 99.8\% of SNI92-to-SNI02 codes and 97.7\% of SNI02-to-SNI07 codes. The joint representation performs somewhat less well on these retrieval diagnostics than the original text representation, again reflecting the trade-off introduced by incorporating additional economic information. Nevertheless, the high cross-version correspondence provides a practical bridge across classification changes and makes it possible to study industries from all three classification periods within a common representation.

The regional evaluation provides the main empirical test of whether this learned geometry contains information relevant to regional development. Opportunity density shows a strong gradient across all three outcomes. Industry-entry and startup-formation rates rise sharply toward the upper density deciles, whereas industry-exit rates decline as density increases. Entry and startup formation are conceptually related---the emergence of an industry can coincide with the founding of new firms---but they capture different events, and the consistency across both outcomes as well as industry exit provides a broader indication that opportunity density reflects conditions associated with the emergence and persistence of regional industrial activity. Across classification periods, the associations are large relative to the corresponding baseline rates and retain the same direction in every reported year: positive for industry entry and startup formation and negative for industry exit. They also persist, with similar magnitudes, when the representation is trained without the municipalities or years in which density is evaluated, and they vary little across training runs.

Multiple regressions with the individual relational views further illustrate the value of their joint representation. Opportunity density retains associations with industry entry, industry exit, and startup formation in almost all models containing a density from one individual view or the untrained text initialization. The exception is industry entry in SNI07 when labour-mobility density is included. Densities based on input--output relatedness or text initialization have confidence intervals including zero once opportunity density is accounted for. Labour mobility is the main complementary view: it remains associated with all three outcomes and is about as strongly associated with industry entry as opportunity density, consistent with the prominent role of skill relatedness in the regional diversification literature~\cite{neffkeSkillRelatednessFirm2013,neffkeHowRegionsDiversify2011}. For industry exit and startup formation, opportunity density shows the larger coefficient magnitudes. These comparisons suggest that the joint representation retains information beyond most individual views, while labour mobility retains information about industry entry that the learned geometry does not fully absorb.

The shared latent representation also provides analytical possibilities beyond the opportunity-density measure used here. Its two-dimensional t-SNE projection offers an exploratory view of industrial neighbourhoods among SNI07 industries, although all quantitative analyses rely on the original 64-dimensional latent geometry rather than the projection. More importantly, the learned embeddings provide a general similarity structure that can be localized through regional industry portfolios. Opportunity density is one direct way of doing so, but future work could investigate other measures derived from the same latent representation to characterize regional opportunity spaces, compare regional industry profiles, and examine how different parts of the industry space are represented within regions.

One limitation of the present framework is that it learns a single, static latent representation of the industry space. Regional industry portfolios evolve over time, but the underlying latent industry geometry itself is held fixed. A natural extension would therefore allow industry embeddings, and thus the latent geometry, to evolve over time. Such an extension was beyond the scope of the present study because the temporal availability and quality of the heterogeneous data sources differ substantially. With more consistent longitudinal information, dynamic representations could distinguish changes in regional industry portfolios from changes in the relational structure among industries themselves.

More broadly, the framework illustrates how machine learning methods---in particular, representation learning---can be used to operationalize latent concepts in economic geography and regional studies. Its architecture is designed not only to integrate heterogeneous relational information, but also to produce an information-fused similarity structure that can support both the analysis of industry relations and their regional evaluation. In this respect, the framework relates to a broader line of work using machine learning to extract meaningful representations from heterogeneous data sources in geographical and regional applications \citep[see also][]{astromEnhancingCarbonEmission2025}. Future extensions could incorporate additional relational views capturing, for instance, market conditions or institutional architectures. Furthermore, by focusing on more recent periods for which richer data are available, temporally evolving latent representations could be developed, allowing co-evolutionary processes with regional opportunity spaces to be analysed. For regional policymakers, such representations could provide a diagnostic basis for exploring how potential activities relate to existing regional industry portfolios and for comparing regional opportunity structures.

\section{Materials and Methods}
\label{sec:methods}

Our framework combines heterogeneous industry relations through multi-view representation learning, then uses the learned geometry to measure opportunity density in regional employment portfolios. Its defining features are separate diffusion targets for each view, one shared cosine geometry across views and classification levels, and log-employment-weighted regional localization. Detailed derivations and implementation settings are provided in SI Appendix, Section~2 (Supporting Methods).

\subsection{Data and relational views}
\label{subsec:source-specific-relations}

We use Swedish employment-register data for 1990--2019 and three versions of the Standard Industrial Classification: SNI92, SNI02, and SNI07~\cite{statisticsswedenscbMIS200722007,europeancommissionNACERev22008}. Four-digit industries are the units of representation; three-digit groups, two-digit divisions, and sections provide their hierarchical context. Each code remains specific to its classification version, while official concordances connect codes across successive versions. This preserves classification differences without treating recoding as economic change.

Labour mobility records annual worker movements between industries within each classification period~\cite{peterfrykblomMetodOchData2021,mondaniEvolvingNetworkLabor2017}. Following the skill-relatedness approach, we compare observed and expected worker flows~\cite{neffkeSkillRelatednessFirm2013}. Expected flows are estimated by Poisson regression using industry employment, growth, and earnings. We add one to observed and expected counts before taking their ratio, then average the two directions. Pairs require at least ten pooled moves. Common ownership records industry co-occurrence within enterprises, normalized by the industries' total co-occurrence counts. Each enterprise--year contributes once per industry pair; pairs require three such observations, and enterprise--years spanning more than 30 industries are excluded.

Input--output relatedness uses Swedish two-digit tables from Statistics Sweden and Eurostat for SNI07, and the World Input--Output Database for SNI02~\cite{statisticsswedenInputOutputTables2019,eurostat_suiot,timmerIllustratedUserGuide2015}. Source years are 2008--2016, 2010--2019, and 1995, 2000, and 2005, respectively. Flows are pooled within each source, symmetrized, and divided by the combined inflows and outflows of the two industries. Aggregated source codes are distributed across their constituent division pairs; overlapping SNI07 measures are averaged. Input--output relations retain their two-digit resolution, with no SNI92 view. Classification information supplies two further constraints: bidirectional parent--child hierarchy links and concordances that distribute weight equally across each code's valid counterparts.

\subsection{Multi-scale relational targets}
\label{subsec:diffusion-targets}

Each relational view is represented as a weighted graph and row-normalized to obtain a transition matrix $P^{(\ell)}$. Graph diffusion captures direct and indirect connectivity at different scales~\cite{coifmanDiffusionMaps2006,gasteigerDiffusionImprovesGraph2019}. We use powers of the transition matrix as scale-specific relational targets:
\begin{equation}
T^{(\ell,s)}=\left(P^{(\ell)}\right)^s.
\label{eq:layer-diffusion}
\end{equation}
Direct links correspond to $s=1$; higher scales incorporate paths through intermediate industries. Labour mobility, common ownership, and hierarchy contribute to direct and higher-order structure. Input--output and classification concordances contribute as direct relations. Diffusion is performed within each view, and its scale-specific targets remain separate during training. Thus, multi-view data fusion occurs through joint learning rather than by averaging the input graphs. The resulting matrices serve as training targets; we do not extract diffusion-map coordinates.

The direct hierarchy target is an ordinal constraint on SNI07 four-digit industries. It orders mean similarities from the same three-digit group, through the same division and section, to different sections. A margin loss penalizes violations between adjacent tiers. This target preserves hierarchy as an ordering rather than treating it as a probability distribution; higher-scale hierarchy targets remain categorical.

\subsection{Learning a shared latent geometry}
\label{subsec:learning-opportunity-space}

Official industry descriptions provide semantic initialization through the frozen \texttt{all-MiniLM-L6-v2} sentence encoder~\cite{wangMiniLMDeepSelfAttention2020} using the Sentence-BERT approach~\cite{reimersSentenceBERTSentenceEmbeddings2019}, following our earlier use of pretrained transformers on text derived from Swedish register data~\cite{Stark2026TextLifeTrajectories}. We reduce its 384-dimensional outputs to 64 dimensions via principal component analysis~\cite{jolliffePrincipalComponentAnalysis2025}. Training then adjusts the four-digit industry embeddings directly. Higher-level codes are represented by the flat mean of their four-digit descendants within the same classification version. Consequently, relations observed at different classification levels constrain the same underlying industry embeddings.

The embeddings define a shared latent similarity matrix $S_{ij}(Z)=\cos(\widetilde z_i,\widetilde z_j)$, including these derived higher-level representations. Each target uses the relevant part of this matrix, with no view-specific projection. For categorical targets, a row-wise softmax converts cosine scores into latent similarity distributions using a learned target-specific temperature. Self-comparisons are excluded, target rows are normalized over the same support, and rows without target mass are omitted. Labour mobility, common ownership, and input--output targets use Kullback--Leibler divergence~\cite{kullbackInformationSufficiency1951}; categorical hierarchy and classification-concordance targets use cross-entropy. The ordinal hierarchy loss acts directly on temperature-scaled cosine scores, without softmax.

All targets jointly refine the same latent geometry. At training epoch $e$, the objective is
\begin{equation}
\mathcal{L}_e=\sum_q w_{q,e}\max\!\left(\ell_{q,e}-h_q,0\right),
\label{eq:joint-loss}
\end{equation}
where $\ell_{q,e}$ is the current target loss, $w_{q,e}$ its weight, and $h_q$ its saturation threshold. Initial weights allocate a budget within each view family using losses relative to an equal-score reference. This prevents adding diffusion scales from automatically increasing a family's initial weight. Moving-average losses subsequently adjust the weights, while saturation stops targets below their threshold from contributing further gradients.

Backpropagation~\cite{rumelhartLearningRepresentationsBackpropagating1986} differentiates the joint objective through the cosine similarities and target-specific losses in one backward pass. Adam~\cite{kingmaAdamMethodStochastic2014} uses the resulting gradients to update the industry embeddings and target temperatures together. Gradients through higher-level means reach their constituent four-digit embeddings. The relational targets, text encoder, and aggregation rule remain fixed, while weighting statistics are updated outside backpropagation. Numerical settings and preprocessing details are specified in the accompanying code repository at \url{https://github.com/PhilippStark/LatentIndustryRepresentations}.

\subsection{Evaluation of the representation}
\label{subsec:embedding-evaluation}

We compare relational views at the common SNI07 two-digit level using Spearman rank correlation, Jaccard overlap of the ten nearest neighbours, centered kernel alignment~\cite{cortesAlgorithmsLearningKernels2012}, and nearest-neighbour mutual-information estimates~\cite{kraskovEstimatingMutualInformation2004}. The latter two summarize matrix alignment and statistical dependence, respectively. Missing pairs are zero-filled for kernel alignment.

We assess reconstruction using Spearman correlations between learned cosine similarities and the observed relations. Input--output reconstruction is evaluated for SNI07 using two-digit hub similarities and, separately, four-digit similarities averaged to division pairs. Hierarchy preservation is summarized by the Spearman correlation between cosine distance and classification-tree distance (1 for the same three-digit group to 4 for different sections). Because 82\% of all SNI07 pairs lie in different sections, the correlation is computed on pairs balanced across the four tiers: all 865 same-group and 3,594 same-division pairs, plus 5,000 pairs drawn at random with a fixed seed from each of the same-section and different-section tiers. Mean similarities by tier use all pairs. Separation of SNI07 sections is summarized by the cosine-distance silhouette score~\cite{rousseeuwSilhouettesGraphicalAid1987}. Cross-version retrieval ranks candidate successors by cosine similarity and reports mean reciprocal rank, mean average precision, and Hit@$k$. The last measures whether at least one valid successor appears among the first $k$ retrieved codes. The baseline is the original 384-dimensional text encoder, before PCA; view-specific evaluation KL uses a fixed temperature distinct from training temperatures (SI Appendix, Table~S1).

These are transductive diagnostics because the underlying relations also guide training. The t-SNE projection~\cite{JMLR:v9:vandermaaten08a} serves only for visualization; quantitative analyses use the original latent geometry. Regional outcomes provide a further assessment of empirical relevance, but their observation periods overlap some training inputs. Spatial and temporal hold-out checks therefore re-estimate the regional models with representations trained without the evaluated municipalities or years (Section~\ref{subsubsec:regional_regression}). The regional models support association claims rather than causal claims.

\subsection{Opportunity density and regional outcomes}
\label{subsubsec:regional_opportunity_density}

Relatedness-density approaches connect an activity to the activities already present in a location~\cite{hidalgoProductSpaceConditions2007,neffkeHowRegionsDiversify2011}. Here, we combine learned cosine similarities with log-employment weights. Let $E_{mit}$ denote employment in industry $i$, municipality $m$, and year $t$. We localize the learned geometry through opportunity density for a focal industry $j$:
\begin{equation}
D_{mjt}=\frac{\sum_{i\ne j}\log(1+E_{mit})\cos(z_i,z_j)}{\sum_{i\ne j}\log(1+E_{mit})}.
\label{eq:opportunity-density}
\end{equation}
The focal industry is excluded, and zero total weight returns zero. Log-employment weights reduce the dominance of the largest local industries. The density profile across industries describes the regional opportunity structure, restricted to its industrial-structural component.

The regional panel aggregates workplace employment by municipality, four-digit industry, and year, with municipal boundaries harmonized over time. We relate density in year $t$ to outcomes in $t+1$, excluding transitions across classification changes. Explanatory years are 1990--2000 for SNI92, 2002--2005 for SNI02, and 2007--2018 for SNI07.

Industry entry is evaluated among industries absent in year $t$, and occurs when meaningful employment appears in $t+1$. Exit is evaluated among industries present in $t$, and counts disappearance only when their preceding activity was meaningful. Meaningful activity requires positive employment whose inverse hyperbolic sine lies within two population standard deviations of the municipality--year mean across active industries. The band requires at least two active industries and positive variation. Our identification of new firms draws on Swedish register studies~\cite{anderssonCharacteristicsPerformanceNew2013,grillitschDoesTimingIntegrating2021}. We evaluate startup formation among previously absent industries, requiring a newly founded, single-establishment firm with at most ten employees and excluding mergers and continuations. It measures firm founding rather than subsequent startup growth.

Employment and wage growth are annual changes in the inverse hyperbolic sine of employment and of mean labour income, respectively. Employment-growth observations require employment in either year; wage-growth observations require more than five workers in both years. SI Appendix, Figs.~S1 and S2 use binary indicators of positive growth and the main opportunity-density specification.

\subsection{Regional estimation and comparisons}
\label{subsubsec:regional_regression}

For each outcome and classification period, we estimate linear probability models
\begin{equation}
Y_{mj,t+1}=\beta D^*_{mjt}+\alpha_{mt}+\gamma_j+\varepsilon_{mj,t+1},
\label{eq:regional-main}
\end{equation}
where $D^*_{mjt}$ is opportunity density divided by its standard deviation over the outcome-specific candidate panel pooled across classification periods. The terms $\alpha_{mt}$ and $\gamma_j$ are municipality$\times$year and industry fixed effects. Municipality$\times$year effects absorb variation common to a municipality in a given year, such as total employment, the number of active industries, and the local business cycle. The coefficient is identified from differences across industries within the same municipality and year, net of persistent industry differences. Because the candidate sets are unbalanced, the fixed effects are absorbed by alternating projections iterated to convergence, after removing singleton observations.

Opportunity density depends on each industry's position in the latent space, motivating allowance for residual dependence across municipalities for the same industry as well as within municipalities. Standard errors are clustered by municipality and industry~\cite{cameronRobustInferenceMultiway2011}. Confidence intervals and $p$-values use a $t$ distribution with $\min(G_{\mathrm{municipality}},G_{\mathrm{industry}})-1$ degrees of freedom, where $G$ denotes the number of clusters; $p$-values below 0.001 are reported as $p<0.001$. Year-specific coefficients are estimated in the same model by interacting density with base-year indicators. Binary-outcome coefficients are reported in percentage points, alongside 95\% confidence intervals and ratios to observed baseline rates. These ratios are scale comparisons, not relative risks. Raw density-decile rates are descriptive.

Robustness checks replace log-employment weights with raw employment and compare arithmetic with geometric aggregation (SI Appendix, Table~S2). Geometric measures first map cosine similarities to the unit interval. We also re-estimate the entry, exit, and startup-formation models with additive municipality, year, and industry fixed effects and with municipality$\times$industry and year fixed effects (SI Appendix, Table~S8 and Fig.~S3). The latter specification yields intervals including zero for all three outcomes in the short SNI02 period.

Hold-out checks retrain the representation without the municipalities or years in which density is evaluated, keeping all other settings. For the spatial hold-out, the 290 municipalities are divided into five folds, stratified by employment size. For each fold, labour mobility and common ownership are rebuilt without the person-years employed there, and the representation is retrained; density in each municipality uses the representation that excluded it. Each retrained representation implies its own density level for every industry, so these pooled models include industry$\times$fold instead of industry fixed effects. For the temporal hold-out, labour mobility, common ownership, and input--output relatedness are restricted to register years up to 2015, 2016, or 2017, and density in the following year is related to outcomes one year later, with industry$\times$year fixed effects. Training stability is assessed with ten representations trained from randomly perturbed initializations (SI Appendix).

To assess whether opportunity density carries information beyond the individual relational views, we construct a density for each view with Eq.~\eqref{eq:opportunity-density}, replacing $\cos(z_i,z_j)$ with the view's relation between industries $i$ and $j$: the row-normalized labour-mobility and common-ownership relations used in training; the input--output relatedness of the industries' two-digit divisions (zero within a division; unavailable for SNI92); the classification-tree similarity (3 for the same three-digit group, 2 for the same division but different groups, 1 for the same section but different divisions, 0 otherwise); and the cosine similarity of the untrained text initialization (PCA, 64 dimensions). For each outcome, period, and view we estimate
\begin{equation}
\begin{split}
Y_{mj,t+1}={}&\beta_O D^{O*}_{mjt}+\beta_V D^{V*}_{mjt}\\
&+\alpha_{mt}+\gamma_j+\varepsilon_{mj,t+1},
\end{split}
\label{eq:regional-multiple}
\end{equation}
with opportunity density $D^{O*}$ and the view's density $D^{V*}$ each divided by its own standard deviation within the classification period. The fixed effects, clustering, and candidate panels are the same as above. The coefficient $\beta_O$ measures the association of opportunity density holding the view's density constant, and vice versa. Within-period scaling differs from the pooled scaling in the main models, so coefficient magnitudes across these two analyses are not directly comparable. Comparisons between the two densities within a multiple regression are also conditional on their respective standard deviations.

\acknow{This work was supported by FORMAS (grant no. 2022-153657) and seed funding from the AI Lund Initiative at Lund University through the project ``Temporal Harmonization of Industry Classifications with Swedish Register Data.'' The work of A.S. was partially supported by eSSENCE (grant no. 156938) and the Swedish National Space Agency (grant no. 2025-00016). The computations were enabled by resources provided by the National Academic Infrastructure for Supercomputing in Sweden (NAISS), partially funded by the Swedish Research Council through grant agreement no. 2022-06725. and LUNARC, the Centre for Scientific and Technical Computing at Lund University (using COSMOS).}

\showacknow{}

\dataavail{Data supporting the analyses, including the evaluation datasets, learned industry embeddings, and matrices representing the relational views, are available at \url{https://github.com/PhilippStark/LatentIndustryRepresentations}. The repository contains all analysis and implementation code and accompanying documentation. The underlying Swedish register data were provided by Statistics Sweden (SCB) for the project ‘Regional tillväxt mot alla odds’ under agreement no. 260875/800 382-4. Access is restricted to authorized research use; the authors cannot redistribute the record-level data. Researchers seeking access must apply directly to SCB. SCB authorized disclosure for this specified research purpose following its confidentiality assessment, which considers risks to individuals and compliance with applicable data-protection and ethical-review requirements.}

\bibliography{references}

\end{document}
